\documentclass[a4paper]{article}

\usepackage{graphicx}
\usepackage{amsmath,amssymb}
\usepackage{minted}
\usepackage{multirow}
\usepackage{float}
\usepackage{todonotes}
\usepackage{forest}
\usepackage{xstring}
\usepackage{xcolor}
\usepackage[most]{tcolorbox}
\usepackage{xparse}
\usepackage{natbib}

\usetikzlibrary{graphs,graphdrawing}
\usegdlibrary{layered}

% for external graphviz
\input{/home/donkarlo/Dropbox/repo/nd_ai_project/src/nd_ai/knoweldge_representation/language/formal/computer/programming/tex/latex/code_clips/external_graphviz.tex}

% for tags
\input{/home/donkarlo/Dropbox/repo/nd_ai_project/src/nd_ai/knoweldge_representation/language/formal/computer/programming/tex/latex/parts/control_sequence/kinds/semtag/semtag.tex}

% for links
\input{/home/donkarlo/Dropbox/repo/nd_ai_project/src/nd_ai/knoweldge_representation/language/formal/computer/programming/tex/latex/code_clips/seehere.tex}

\title{Planning}
\date{}

\begin{document}
    \maketitle
    \todo{Planning is not decision making, planning uses decion making to choose the best action at each step}
    \todo{Biologically , it is called action planning}
    
    Planning is the process of thinking regarding the activities required to achieve a desired goal. Planning is based on foresight, the fundamental capacity for mental time travel. \seeherefooter{https://en.wikipedia.org/wiki/Planning}
    
    
    \section{From neorology perpective}
        Planning is one of the executive functions of the brain, encompassing the neurological processes involved in the formulation, evaluation and selection of a sequence of thoughts and actions to achieve a desired goal
        \seeherefooter{https://en.wikipedia.org/wiki/Planning#Neurology}
    
    
    \section{Cognitive planning}
        
        Cognitive planning is one of the executive functions. It encompasses the neurological processes involved in the formulation, evaluation and selection of a sequence of thoughts and actions to achieve a desired goal.
        \seeherefooter{https://en.wikipedia.org/wiki/Planning_(cognitive)}
    
    
    \section{Automated planning and scheduling}
        is a branch of artificial intelligence that concerns the realization of strategies or action sequences, typically for execution by intelligent agents \seehere{https://en.wikipedia.org/wiki/Automated_planning_and_scheduling}
        
        Planning is the process of thinking regarding the activities required to achieve a desired goal. \seehere{https://en.wikipedia.org/wiki/Planning}
    
    
    \section{Motor planning}
        \clickurlfooter{https://en.wikipedia.org/wiki/Motor_planning} defines motor planning is a set of processes related to the preparation of a movement that occurs during the reaction time (the time between the presentation of a stimulus to a person and that person's initiation of a motor response).
    
    
    \section{Motion planning}
        \clickurlfooter{https://en.wikipedia.org/wiki/Motion_planning} defines motion planning, also path planning (also known as the navigation problem or the piano mover's problem) as a computational problem to find a sequence of valid configurations that moves the object from the source to destination.
        
        
        \bibliographystyle{plainnat}
        \bibliography{/home/donkarlo/Dropbox/repo/nd_shared_research/bibliography/bibliography.bib}
\end{document}